\documentclass[11pt]{article}
\usepackage[margin=1in]{geometry}
\usepackage{amsmath}
\usepackage{amssymb}
\usepackage{hyperref}
\usepackage{enumitem}

\title{MAF/MADE Residual and Masking Notes (TTZ)}
\author{}
\date{\today}

\begin{document}
\maketitle

\section{Scope of This Note}
This note is only about the architecture discussion: how masking works in MADE,
how residual/skip layers are implemented in the TTZ flow model, and how these
pieces combine inside MAF/MAFMADEMOG.

The active TTZ settings discussed were:
\begin{itemize}[leftmargin=*]
\item \texttt{num\_flows = 5}
\item \texttt{num\_hidden\_layers = 4}
\item \texttt{hidden\_layer\_dim = 384}
\item \texttt{res\_layers\_in\_block = 2}
\item \texttt{conv1x1 = True}
\item \texttt{maf\_residuals = True}
\end{itemize}

\section{Big Picture: What Problem MAF Solves}
Normalizing flows learn an invertible map between data $x$ and latent $u$ with
tractable log-density:
\[
\log p_X(x) = \log p_U(u) + \log\left|\det\frac{\partial u}{\partial x}\right|.
\]
In MAF, each block is autoregressive, so the Jacobian is triangular and
its log-determinant is cheap to compute.

The model gains flexibility by stacking multiple autoregressive transforms,
interleaved with invertible mixing/permutation layers.

\section{How Autoregressive Masking Works in MADE}
\subsection*{Core rule}
Each output dimension $i$ is only allowed to depend on previous dimensions in a
chosen order. In this codebase the ordering is sequential by default.

The masking function constructs binary masks so that illegal connections are
exactly zeroed out in \texttt{MaskedLinear}:
\[
\text{out} = W \odot M \; x + b.
\]
So masking is not a soft regularizer. It is a hard structural constraint.

\subsection*{Degree-based construction}
The implementation assigns each neuron a degree $m$ and allows connections only
from earlier/equal degrees according to the MADE rules:
\begin{itemize}[leftmargin=*]
\item hidden-layer masks use a \emph{non-strict} condition ($\ge$),
\item output mask uses a \emph{strict} condition ($>$).
\end{itemize}

That strict final mask is what guarantees output $i$ cannot use its own input
$x_i$ directly in the autoregressive parameter prediction.

\section{Gaussian Parameterization Used in MAF Blocks}
The Gaussian MADE variants output two vectors per feature:
\[
\mu(x_{<i}),\quad \alpha(x_{<i})
\]
where $\alpha$ is log-scale. The transform used in normalizing direction is:
\[
u_i = (x_i - \mu_i)\exp(-\alpha_i).
\]
Then:
\[
\log\left|\det\frac{\partial u}{\partial x}\right| = -\sum_i \alpha_i.
\]

Sampling (inverse direction) is sequential in feature index because each
$x_i$ needs already-generated previous coordinates.

\section{Residual Blocks: What ``Skipping'' Means Here}
\subsection*{Form of the residual update}
Inside \texttt{ResMaskedLinear}, the block computes:
\[
y = x + f(x),
\]
where $f$ is a masked stack of nonlinear layers.

This is the skip connection we discussed. The identity path ($x$) is always
present, and the block learns a correction $f(x)$ instead of a full remapping.

\subsection*{Why this helps}
Residual form improves optimization because:
\begin{itemize}[leftmargin=*]
\item gradients can flow through the identity branch,
\item deep stacks are less prone to representational collapse,
\item near-identity behavior is easy early in training.
\end{itemize}

\subsection*{Important implementation detail: zero initialization}
The final layer inside each residual mini-network is zero-initialized in
\texttt{ResMaskedLinear} (weights and bias).

Consequence at initialization:
\[
f(x) \approx 0 \quad\Rightarrow\quad y \approx x.
\]
So each residual block starts near identity and gradually learns structured
deviations. This typically stabilizes early training dynamics.

\section{Two Depth Knobs You Asked About}
In \texttt{GaussianResMADE(input\_size, k, n\_blocks, l)}:
\begin{itemize}[leftmargin=*]
\item \texttt{n\_blocks} is the number of residual blocks.
\item \texttt{l} is the number of masked layers inside each residual block.
\end{itemize}

Mapping to TTZ config in this repository:
\begin{itemize}[leftmargin=*]
\item \texttt{n\_blocks = num\_hidden\_layers = 4}
\item \texttt{l = res\_layers\_in\_block = 2}
\item \texttt{k = hidden\_layer\_dim = 384}
\end{itemize}

So one autoregressive block contains 4 residual units, each with 2 internal
masked layers, all at width 384.

\section{How One Full MAFMADEMOG Block Is Built}
For each of the \texttt{num\_flows} repeats, the code adds:
\begin{enumerate}[leftmargin=*]
\item optional \texttt{BatchNormFlow},
\item either \texttt{Conv1x1PLU} (if \texttt{conv1x1=True}) or \texttt{ReverseFlow},
\item an autoregressive transform: \texttt{GaussianResMADE} if residuals are enabled.
\end{enumerate}

With your settings, this means each repeat uses batch norm, then invertible
1x1 mixing, then residual autoregression.

\section{Why 1x1 Invertible Mixing Is Useful with Autoregression}
Autoregressive layers are order-sensitive. Without mixing/permutation, later
blocks repeatedly see very similar coordinate ordering constraints.

The invertible 1x1 linear transform rotates/mixes channels between blocks,
which changes what each next autoregressive layer sees as ``earlier'' vs
``later'' information content.

Intuition: masking gives triangular structure; 1x1 mixing prevents that
structure from becoming too rigid across depth.

\section{Forward vs Inverse Cost Asymmetry (Important Practical Point)}
In autoregressive flows:
\begin{itemize}[leftmargin=*]
\item density evaluation in normalizing direction is parallel over dimensions,
\item sampling in inverse direction is sequential over dimensions.
\end{itemize}

This is why MAF is typically efficient for likelihood training and comparatively
slower for generation than coupling-based alternatives.

\section{How Skip Connections and Masks Interact}
A common confusion is whether residual skips break autoregressive validity.
They do not, because:
\begin{itemize}[leftmargin=*]
\item the residual branch itself uses masked linear layers,
\item the identity add $x + f(x)$ is within hidden feature space of fixed width,
\item the output projection still uses the autoregressive output mask.
\end{itemize}

So residual design improves optimization capacity while preserving the same
autoregressive dependency constraints.

\section{Concrete Parameter Counting Intuition (Qualitative)}
Increasing these knobs does the following:
\begin{itemize}[leftmargin=*]
\item \texttt{hidden\_layer\_dim} increases per-layer capacity strongly.
\item \texttt{num\_hidden\_layers} increases number of residual units per MAF block.
\item \texttt{res\_layers\_in\_block} deepens each residual unit.
\item \texttt{num\_flows} repeats the whole (mixing + autoregressive) macro-block.
\end{itemize}

In practice, these knobs trade off expressivity, train time, memory use, and
stability. Residual and near-identity initialization are key enablers when
depth is increased.

\section{Practical Reading of Your Current Setup}
Given your TTZ settings (5 flow repeats, residual MADE blocks, width 384,
4 residual blocks each with 2 internal masked layers, and invertible 1x1 mixing):
\begin{itemize}[leftmargin=*]
\item the model has substantial capacity for nontrivial correlations,
\item residualization should help keep training stable at this depth,
\item 1x1 mixing should reduce repeated-order artifacts across flow repeats.
\end{itemize}

This is exactly the architecture pattern one would choose when simple
autoregressive stacks are underfitting correlation structure.

\section{Worked Example: One Event Through One Flow Repeat}
This section gives a concrete step-by-step pass for a single event vector
$x\in\mathbb{R}^{15}$ (TTZ input dimension in the current setup).

Assume we are in one repeat of the flow stack with:
\begin{enumerate}[leftmargin=*]
\item \texttt{BatchNormFlow},
\item \texttt{Conv1x1PLU},
\item \texttt{GaussianResMADE}.
\end{enumerate}

\subsection*{Step 1: Batch normalization transform}
Input event is $x^{(0)}=x$.
Batch norm applies a per-feature affine normalization:
\[
x^{(1)}_j = w_j\frac{x^{(0)}_j-\mu_j}{\sqrt{\sigma_j^2+\epsilon}} + b_j.
\]
This contributes an additive log-Jacobian term
$\log|\det J_{\text{BN}}|$ that is easy to compute from per-feature scales.

Interpretation: this step re-centers/re-scales channels so later layers operate
in a numerically friendlier regime.

\subsection*{Step 2: Invertible 1x1 linear mixing}
Then \texttt{Conv1x1PLU} applies an invertible linear map:
\[
x^{(2)} = x^{(1)}W,
\]
where $W$ is parameterized through PLU factors for stable invertibility and
cheap log-determinant:
\[
\log|\det J_{\text{mix}}| = \log|\det W|.
\]

Interpretation: coordinates are rotated/mixed so the next autoregressive module
does not always see the same axis-aligned dependency pattern.

\subsection*{Step 3: Residual autoregressive map (GaussianResMADE)}
Now \texttt{GaussianResMADE} receives $x^{(2)}$ and predicts two 15-vectors:
\[
\mu(x^{(2)}_{<i}),\quad \alpha(x^{(2)}_{<i}).
\]

Each hidden computation inside this network uses masked linear layers. Residual
blocks apply hidden-state updates of the form:
\[
h \leftarrow h + f(h),
\]
with masked $f$. Because the final masked layer in each residual unit is
zero-initialized, each unit starts near identity at initialization.

The autoregressive transform for each feature is:
\[
u_i = (x^{(2)}_i - \mu_i)\exp(-\alpha_i).
\]
Collecting all features gives $u^{(3)}$ and log-Jacobian term:
\[
\log|\det J_{\text{AR}}| = -\sum_i \alpha_i.
\]

\subsection*{Step 4: Total contribution of this repeat}
For this one repeat, output is $u^{(3)}$, and the repeat contributes:
\[
\log|\det J_{\text{repeat}}|=
\log|\det J_{\text{BN}}|+
\log|\det J_{\text{mix}}|+
\log|\det J_{\text{AR}}|.
\]

This output is then fed to the next repeat (out of 5 total in your config).

\subsection*{What changes during sampling (inverse direction)}
During generation, the inverse order is used inside each repeat:
\begin{enumerate}[leftmargin=*]
\item inverse autoregressive step (sequential in feature index),
\item inverse 1x1 mixing,
\item inverse batch norm.
\end{enumerate}

The sequential part is specifically the autoregressive inverse:
\[
x_i = u_i\exp(\alpha_i) + \mu_i,
\]
where $\mu_i,\alpha_i$ depend on already reconstructed earlier coordinates.

This is the core reason sampling is slower than density evaluation in MAF.

\section{Minimal Mental Model}
If you want one compact memory aid:

\begin{quote}
MADE masks decide \emph{who can talk to whom}. Residual skips decide
\emph{how hard optimization has to work}. 1x1 invertible mixing decides
\emph{whether each new autoregressive block sees a fresh coordinate geometry}.
\end{quote}

Together, they give an expressive but still tractable flow for high-dimensional
physics features.

\end{document}
